\topsection[sec:project]{Project management}

\subsection{Plan and milestones}
The scoping note planned four milestones over the eight weeks, with an end-to-end chain first so that
integration would never be left to the end. Table~\ref{tab:plan} compares the plan with the dated evidence.

{\tablefont
\begin{xltabular}{\textwidth}{P{4.4cm}P{1.9cm}L}
\caption{Milestones planned in the scoping note and when each one was reached.}\label{tab:plan}\\
\toprule
\headrow\hd{Milestone} & \hd{Planned} & \hd{Reached} \\
\midrule
\endfirsthead
\toprule
\headrow\hd{Milestone} & \hd{Planned} & \hd{Reached} \\
\midrule
\endhead
\bottomrule
\endfoot
Minimum chain: register, hash, sign, verify, detect tampering & Week 2 & 5~October: version 0.2 already attested the four stages, with 113 tests and 26 attack and control variants in five scenarios. First progress report \\
Four-stage chain with roles, policy and revocation & Week 4 & 16~October: version 0.3 added the revocation experiment and external models, with 143 tests and 44 variants. The Compose stack ran on a team laptop. Second progress report \\
Attack scenarios and measurements & Weeks 5 to 7 & 26 variants on 5~October, 44 on 16~October, 48 in the final run of 9~November, each time with the cost measurements \\
Report, paper and demonstration & Week 8 & 9 and 10~November: version 0.4 with the signed trust log and the export, full validation run, public repository with both CI workflows passing, the paper, the website and this report \\
\end{xltabular}}

Each full run is kept: the results of 5~October and 16~October are archived unchanged in the repository
(\repodir{results}), and the experimentation journal has one dated entry per run, which is never deleted.

\subsection{Team and ownership}
Table~\ref{tab:team} gives the area each member is responsible for. It describes ownership, not a record of
individual contributions.

\begin{table}[H]
\caption{Team roles and areas of ownership.}
\label{tab:team}
\tablefont
\begin{tabularx}{\textwidth}{P{3.1cm}P{3.5cm}L}
\toprule
\headrow\hd{Member} & \hd{Role} & \hd{Owns} \\
\midrule
Karan DESHMUKH & Technical lead, security architecture & Architecture and integration, verifier and trust-log design, external-model design, deployment gateway, API, interoperability export, coordination, paper \\
JOGI Sudheer reddy & AI and ML engineer & Dataset and training pipeline, backdoor experiments (experiments 6 and 7), inference benchmark \\
DAMA Dola Krishna & Testing, documentation and quality & Test suite, experimentation journal, collection of results, validation script, dashboard support, reports \\
MEKALA Rambabu & Security and DevOps engineer & Hashing, signatures, key and publisher-key handling, revocation and trust-log operations, Docker, CI/CD \\
\bottomrule
\end{tabularx}
\end{table}

\subsection{Reporting to the referent}
The brief asks for a detailed weekly progress report. Karan DESHMUKH sent it to the referent by email. Two
formal progress reports mark the main steps: the report on the initial period (5~October 2026) and Progress
Report~2 (16~October 2026). The progress update for version 0.4 (9~November 2026) is in the repository
(\repofile{docs/16_progress_update_v0.4.md}). The two formal reports are archived with this report so that the
progression can be checked.

\subsection{Method}
\textbf{Evidence rule.} From the first report, a feature is reported as completed only when code and
repeatable execution evidence exist, and every number is tied to the run that produced it. This report follows
the same rule.

\textbf{Use of AI assistance.} As stated in each progress report, the code was produced with AI-assisted code
generation and validated by execution: tests, complete validation runs in the reference environment, the CI
workflows and the Docker run on a team laptop. The design decisions, the experiments, their analysis and the
conclusions are the responsibility of the team.

\subsection{Difficulties and lessons}
\begin{itemize}
\item \textbf{Our own experiment broke our design.} In version 0.3 the revocation experiment showed that
  clearing one flag in the database undid a revocation. We reported it as a limitation on 16~October and
  removed it in version 0.4 with the signed trust log. Lesson: attack your own trust anchors, and write the
  limitation down before fixing it.
\item \textbf{A detector must be measured on clean data first.} The first version of the corner-trigger check
  flagged 11 clean digits. Requiring the minimum of the four pixels instead of their mean brought the false
  alarms to zero, on the reference split and on the ten other seeds of Section~\ref{sec:seeds}.
\item \textbf{Run the automation before trusting it.} The CI guard against committed private keys matched
  its own search pattern in the workflow file and would have failed on every run. Executing the workflow steps
  locally found the bug before the first run on GitHub.
\item \textbf{One environment is not enough.} Docker could not run in the reference environment, so the
  Compose stack was executed on a team laptop. Timings still come from one environment.
\item \textbf{Saying what a result does not mean is part of the result.} Keeping \st{VERIFIED} from being read
  as ``safe'' shaped the statuses, the dashboard and the wording of every report.
\end{itemize}
